\subsection{Expansion of the collected analytic formula}
\label{sec:comp-total-energy-low}
The complete product form \eqref{eq:comp-total-energy-product} gives
two further orders beyond the leading result above. The external energy variable remains $X$.
Direct expansion gives
\begin{equation}
\begin{aligned}
 \sigma_1=\frac{\alpha_{\mathrm T}^3}{m^2}\biggl\{
 &X^2\left[-\frac{16}{9}\ln X+\frac7{15}\right]
 +X^3\left[\frac43\ln X-\frac{113}{90}\right]\\
 &+X^4\left[-\frac{25}{9}\ln X+\frac{203}{125}\right]
 +O\!\left(X^5|\ln X|\right)\biggr\},\qquad X\to0^+.
\end{aligned}
\label{eq:comp-collected-total-low-four}
\end{equation}
No individual virtual or real integral is expanded again here. The
starting point is the complete infrared-finite analytic formula.

The expansion can be performed recursively from the defining GPL
kernels. For a nonzero rational letter $a$, and for the root letter,
\begin{align}
 \frac1{X-a}&=-\sum_{j=0}^{\infty}\frac{X^j}{a^{j+1}},
                      \qquad |X|<|a|,\\
 \frac1{\sqrt{X(X+4)}}
 &=\frac1{2\sqrt X}\sum_{j=0}^{\infty}
       \binom{-1/2}{j}\frac{X^j}{4^j},\qquad |X|<4.
\end{align}
For the zero letter the kernel is $1/X$. Starting with the empty
word $G(\varnothing;X)=1$, multiply the kernel series by the series
for the shorter word and integrate each term. For $p\ne-1$,
\begin{equation}
 \int X^p\ln^k X\,\dd X
 =X^{p+1}\sum_{j=0}^{k}
       \frac{(-1)^j k!}{(k-j)!(p+1)^{j+1}}\ln^{k-j}X.
\end{equation}
For $p=-1$ the primitive is $\ln^{k+1}X/(k+1)$.
The zero-letter words are normalized as
$G(\underbrace{0,\ldots,0}_{k};X)=\ln^k X/k!$.
All other integration constants follow from their vanishing limit at
the lower endpoint. This gives the required half-integer power series
without a numerical fit.

The rational coefficients of the total have poles of order at most
four at $X=0$. Accordingly the GPL series are computed through $X^8$
before multiplication, which suffices for the total through $X^4$.
On summing all terms, the negative powers, the constant and linear
terms, every half-integer power, and powers of $\ln X$ above one
cancel exactly. The first surviving term is the $X^2$ term displayed
in \eqref{eq:comp-collected-total-low-four}. The local kernel series
converge in a sufficiently small interval $0<X<1$.
After the displayed terms are removed, the next possible contributions
have the form $X^5$ times an analytic function or $X^5\ln X$ times an
analytic function, giving the stated remainder. Indeed, every
non-elementary word in the product form has at most one trailing zero,
and the only explicit factor $\ln X$ multiplying a GPL multiplies a
word with no trailing zero. The elementary part is also at most linear
in $\ln X$. Thus no higher logarithmic power can first appear in the
omitted orders.

Only after deriving these coefficients do we compare the leading
term with Lee, Schwartz and Zhang~\cite{qed-processes-lsz},
Equation (10). Their variable is the same $X=(s_{\mathrm{phys}}-m^2)/m^2$,
and their order-$\alpha^3$ normalization is $\alpha^3/m^2$.
The logarithmic coefficient $-16/9$ and the finite coefficient $7/15$
agree separately. That displayed literature equation does not give
the $X^3$ and $X^4$ coefficients, so no comparison of those two orders
is inferred from it.
